\subsection{Supuestos del modelado}

\begin{itemize}
  \item La foto contiene la información del BCS. Las vistas posterior y lateral muestran lo suficiente de cadera, base de la cola y costillas para estimar el BCS.
  \item La etiqueta es confiable. El BCS asignado por los expertos del CAETEC se toma como verdad de referencia, y todas las métricas se calculan contra él. La etiqueta de cada vaca es el promedio de sus evaluaciones, respetando el redondeo con incrementos de 0.25.
  \item La vaca es la unidad de análisis. Las fotos de una misma vaca no son independientes; por eso la partición es por vaca, la pérdida pondera a cada vaca por igual y las métricas se calculan por vaca.
  \item El BCS no cambia entre las fotos de una vaca en una sesión. La predicción de la vaca es el promedio (o la mediana) de sus fotos.
  \item El volteo horizontal no cambia la etiqueta. El BCS es simétrico, así que el volteo se usa en aumentación y en la inferencia.
  \item Las características aprendidas con imágenes generales transfieren al ganado, y las aprendidas con fotos de anca de otras granjas transfieren al CAETEC. Los pasos 2 y 3 lo ponen a prueba.
  \item La resolución de entrada ($384\times288$ en la mayoría de los modelos) conserva el contorno corporal relevante.
  \item Quitar el fondo de la foto no elimina información útil.
  \item La escala de reporte es la del experto: 1 a 5 en cuartos de punto.
\end{itemize}
